\documentclass{beamer}
\usetheme{Warsaw}
\usecolortheme{beaver}
\usepackage{amsmath}

\author{Mohammadhossein Farzanfar}
\title{A Brief Introduction to LaTeX for Scientific Writing}
\institute{RUDN University}

\begin{document}

\begin{frame}
\titlepage
\end{frame}

\begin{frame}{Main Message: Why Use LaTeX?}
\begin{block}{Key Idea (Covered Early)}
LaTeX is great for scientific documents because it handles math, references, and structure logically – not WYSIWYG like Word.
\pause
Example: Unlike word processors, LaTeX uses markup for sections.
\end{block}
\end{frame}

\begin{frame}{Math in LaTeX}
A \alert{set} is a collection of objects. \uncover<2->{For example:}
\[ Z = \{\text{cow}, \text{pig}, \text{elephant}\}. \]
\uncover<3->{We call the objects in $Z$ the \alert{elements} of $Z$.}
\uncover<4->{Frequently encountered sets:}
\begin{align*}
\uncover<5->{\mathbb{N}} &= \{1,2,3,\ldots\} \quad (\text{natural numbers}) \\
\uncover<6->{\mathbb{Z}} &= \{\ldots,-2,-1,0,1,2,\ldots\} \quad (\text{integers})
\end{align*}
\end{frame}

\begin{frame}{Example: Derivative (Before Definition)}
\begin{block}{Example First}
The derivative of $f(x) = x^2 \cdot \sin(x)$:
\begin{align*}
f'(x) \uncover<2->{&= 2x \cdot \sin(x) +} \uncover<3->{x^2 \cdot \cos(x).}
\end{align*}
\pause
\end{block}
\begin{block}{Definition}
Product rule: If $f(x) = g(x) \cdot h(x)$, then $f'(x) = g'(x) \cdot h(x) + g(x) \cdot h'(x)$.
\end{block}
\end{frame}

\begin{frame}{Tips for Presentations}
\begin{itemize}
\item<1-> Use pictures over text (e.g., include diagrams).
\item<2-> Give examples before abstract definitions.
\item<3-> Keep it short – time yourself!
\end{itemize}
\end{frame}

\end{document}
